% ---------------------------------------------------------------------
%  2. Mechanism: Lean-gated coordination (Part One design)
% ---------------------------------------------------------------------
\section{Mechanism: Lean-gated coordination}
\label{sec:mechanism}

\subsection{Design principle: Lean at every step, not only at the end}
The comparator judges once. My system uses Lean far more often than that: the
unit of work is a loop, not a single generation.
\[
\text{candidate}
\;\xrightarrow{\text{Lean}}\;
\text{structured diagnostic}
\;\xrightarrow{\text{model}}\;
\text{repair}
\;\xrightarrow{\text{Lean}}\;
\cdots
\]
Every candidate is compiled against a clean Mathlib environment; the service
returns typed messages, a sorry flag, and timing. Those messages --- not a
model's opinion of its own proof --- decide what happens next. The consequence I
rely on throughout: \textbf{cross-model coordination happens over the compiler's
output, never over prose.} When a stuck proof passes from one model to the other,
it carries the failed candidate plus the exact Lean error, and the second model
is asked to repair. There is no natural-language debate between agents, because a
Lean diagnostic is a stronger, cheaper, and more objective signal than one model
critiquing another in English. Three invariants hold at every rung: no sorry or
admit, no new axioms, and no silent rewriting of the theorem statement --- so
that ``Lean accepted it'' means to me what it will mean to the grader.

\subsection{A finite escalation ladder}
The agent is a \emph{finite} ladder. Success returns immediately; on failure the
next rung runs until it stalls, hits its cap, or the ladder is exhausted. The
order is deliberate --- cheapest and most reliable first, most expensive and most
speculative last (Figure~\ref{fig:ladder}):

\begin{enumerate}
  \item \textbf{T0 --- tactic sweep.} Zero-model closing tactics on the locked
        challenge. Free wins on easy goals, no API spend.
  \item \textbf{Same-model repair champions (Qwen, then GPT-OSS).} One model
        proposes, Lean checks, the \emph{same} model repairs from the failed
        proof and the exact messages, under a turn budget.
  \item \textbf{Seeded cross-model handoff.} The best failed Qwen candidate is
        repaired by GPT-OSS, and vice versa --- carrying the Lean context
        forward, not restarting.
  \item \textbf{Independent slot / whole fallback.} Split the file into locked
        declarations, sample and repair each, recombine under a strict integrity
        check. This is where cheap multi-sampling pays off, and where hard
        multi-theorem problems tend to fall.
  \item \textbf{Bounded program/bridge experiment}, then a \textbf{residual
        lemma-banking stage} (\S\ref{sec:frontier}) if wall remains.
  \item \textbf{Stop} --- accept, stall, or exhausted.
\end{enumerate}

The ladder is the design thesis made concrete: it starts from ``less'' (a free
tactic sweep) and escalates to ``more'' only as the problem resists, and it never
averages the two models --- it \emph{routes} between them and lets Lean
adjudicate each step. This routing is exactly what lets one policy collect wins
that are otherwise split across the two models (\S\ref{sec:results}).

\begin{figure}[h]
\centering
\begin{tikzpicture}[
  font=\small,
  box/.style={draw, rounded corners=2pt, align=center, inner sep=4pt,
              minimum height=0.72cm},
  arr/.style={-{Latex}, thick}
]
\node[box, fill=black!6]   (t0)  {T0\\ sweep};
\node[box, right=0.32cm of t0, fill=linknavy!12] (rq) {repair\\ Qwen};
\node[box, right=0.32cm of rq, fill=linknavy!12] (rg) {repair\\ GPT};
\node[box, right=0.32cm of rg, fill=orange!16]   (h)  {handoff\\ Q$\leftrightarrow$G};
\node[box, right=0.32cm of h,  fill=green!14]    (fb) {slot\\ fallback};
\node[box, below=0.55cm of rg, fill=yellow!22]   (exp){program /\\ bridge};
\node[box, right=0.5cm of exp, fill=red!10]      (res){residual\\ banking};
\node[box, right=0.5cm of res, fill=black!5]     (stop){accept /\\ exhausted};
\draw[arr] (t0) -- (rq);
\draw[arr] (rq) -- (rg);
\draw[arr] (rg) -- (h);
\draw[arr] (h)  -- (fb);
\draw[arr] (fb.south) -- ++(0,-0.24cm) -| (exp.north);
\draw[arr] (exp) -- (res);
\draw[arr] (res) -- (stop);
\end{tikzpicture}
\caption{The finite coordination ladder. Lean accepts or rejects at every rung;
cheap and reliable rungs run first. Residual banking is post-ladder spend, not
the spine.}
\label{fig:ladder}
\end{figure}

\subsection{Acceptance hygiene, and why it is a feature}
A proof accepted by a warm REPL during search is not automatically a proof I
submit. Every final candidate passes a \textbf{strict integrity} guard:
declaration signatures are immutable up to the proof body, answer shapes and
numeric values are checked against the locked problem, and known circular tricks
are rejected. This matters more than it sounds. A model can ``solve'' a Putnam
answer-set problem by \emph{redefining the answer set to be the question} --- for
\texttt{putnam\_2018\_a1}, declaring the solution set to be exactly the pairs
satisfying the left-hand side, so the required equivalence collapses to a
reflexivity and the comparator accepts a proof that establishes nothing. My guard
infers the answer shape from the challenge and rejects such
set-builder / $\lambda$-plus-reflexivity tautologies. I therefore report
\emph{substantive} passes only. Everything counted in this report is a proof of
the actual statement --- a stricter and more defensible bar than raw comparator
acceptance, and the reason my headline numbers hold up when they are re-run.
